% Avsnitt 10.8, ur lectures/L10/README.md.

\section{Sammanfattning}
\label{c10:sec:review}

\needspace{6\baselineskip}
\subsection*{Det här ska du kunna nu}
\begin{itemize}
  \item Säga varför kommunikationsprotokoll använder ramar, checksummor och arbitrering.
  \item Beskriva CAN-ramens format fält för fält, och konstruera en komplett ram för hand.
  \item Förklara hur wired-AND-arbitrering löser upp konkurrens utan bussmästare, och avgöra vilken
    av två noder som vinner.
  \item Tillämpa stoppbitsregeln på en bitföljd, för hand.
  \item Skriva \file{can_def.vhd}: konstanterna, subtyperna och CRC-polynomet som resten av
    projektet läser.
  \item Redogöra för gruppens arbetssätt: branchning, PR:er och granskning.
\end{itemize}

\needspace{6\baselineskip}
\subsection*{Frågor att testa dig själv med}
\begin{enumerate}
  \item Varför behöver ett protokoll ett längdfält när ramen ändå har en tydlig början?
  \item Varför har CAN ingen destinationsadress, och vad använder den i stället?
  \item Två noder börjar sända samtidigt. Vilken vinner, och vad gör förloraren?
  \item Varför måste en sändare stoppa in en bit efter fem lika i rad, och varför gäller regeln inte
    hela ramen?
  \item Vad händer med CRC-beräkningen för en instoppad bit?
  \item Varför ligger sampelpunkten runt 70~\% in i bitperioden i stället för i mitten?
\end{enumerate}

\needspace{6\baselineskip}
\subsection*{Riktvärde}
\file{can_def.vhd} bör vara skriven och mergad när passet är slut. Det är en kort fil och resten av
projektet läser den, så den är värd att göra klar direkt.

Riktvärdena i de här kapitlen är just riktvärden. De är inte milstolpar, de betygsätts inte, och
gruppen bestämmer själv i vilken ordning modulerna byggs; se bilaga~\ref{app:project}.

\needspace{6\baselineskip}
\subsection*{Blicken framåt}
\Chapref{c11:ch} går från protokoll till blockschema: kontrollerns arkitektur, den toppnivå varje
senare modul instansieras i, registerkartan som drivrutinsklassen ska konsumera, och projektets
simuleringsflöde i GHDL.
